% Generated draft: requires mathematical and source review.
\documentclass[11pt]{article}
\usepackage[margin=25mm]{geometry}
\usepackage{fontspec}
\setmainfont{Latin Modern Roman}
\usepackage{amsmath,amssymb,mathtools,unicode-math}
\setmathfont{Latin Modern Math}
\usepackage{enumitem,needspace,xurl,hyperref}
\setlength{\parindent}{0pt}
\setlength{\parskip}{5pt}
\setlength{\emergencystretch}{3em}
\newcommand{\N}{\mathbb{N}}
\newcommand{\Z}{\mathbb{Z}}
\newcommand{\Q}{\mathbb{Q}}
\newcommand{\R}{\mathbb{R}}
\newcommand{\C}{\mathbb{C}}
\newcommand{\F}{\mathbb{F}}
\newcommand{\A}{\mathbb{A}}
\newcommand{\PP}{\mathbb{P}}
\newcommand{\E}{\mathbb{E}}
\newcommand{\eps}{\varepsilon}
\newcommand{\dd}{\,\mathrm{d}}
\begin{document}
\section*{UnsolvedMath ID: 801146}
\textbf{Source problem code:} catalogue entry
\subsection{Definitions and mathematical statement}
A \(\mathrm{II}_1\) factor is a von Neumann algebra with center \(\C 1\), no nonzero minimal projections, and a faithful normal tracial state \(\tau_A\) normalized by \(\tau_A(1)=1\). Interpret separability as separability of the predual. Write \(\|x\|_{2,A}=\tau_A(x^*x)^{1/2}\), and let \(\overline{\otimes}\) denote the spatial von Neumann algebra tensor product.

The source uses a common ultrafilter \(\mathcal U\) without specifying its index set or whether it is nonprincipal. For a fixed index set \(I\) and ultrafilter \(\mathcal U\) on \(I\), the tracial ultrapower notation means
\[
A^{\mathcal U}=\ell^\infty(I,A)\big/\{(a_j)_{j\in I}:\lim_{j\to\mathcal U}\|a_j\|_{2,A}=0\},
\]
where \(\ell^\infty(I,A)\) consists of operator-norm bounded families. Write \(A\equiv B\) if \(A\) and \(B\) have the same values for every sentence in the continuous first-order language of tracial von Neumann algebras.

For the common \(\mathcal U\) intended by the source, is the following implication true for every four separable \(\mathrm{II}_1\) factors \(N_1,N_2,M_1,M_2\)?
\[
N_1^{\mathcal U}\overline{\otimes}N_2^{\mathcal U}
\cong
M_1^{\mathcal U}\overline{\otimes}M_2^{\mathcal U}
\quad\Longrightarrow\quad
\exists\sigma\in S_2\ \forall i\in\{1,2\},\quad N_i\equiv M_{\sigma(i)}.
\]
Here \(\cong\) denotes an isomorphism of von Neumann algebras and \(S_2\) is the permutation group of \(\{1,2\}\). The ultrapowers are taken separately before forming the tensor products. The intended restrictions on \(\mathcal U\), and the scope of its quantification, remain unspecified in the supplied record.
\subsection{Short English statement}
If tensor products of the respective tracial ultrapowers of two pairs of separable \(\mathrm{II}_1\) factors are isomorphic, must the original factors have matching elementary theories after possibly interchanging one pair? The source labels this question open; its intended ultrafilter assumptions require clarification.
\subsection{Sources}
\begin{itemize}
\item Local UnsolvedMath registry, permanent ID \(801146\), source problem code \texttt{catalogue entry}, title ``Ultraproduct \(\mathrm{II}_1\) Factors --- Problem 48.'' The supplied record attributes the problem to Elayavalli, gives the proposed year as 2025, and labels its status \emph{open}. This status is supplied curation, not an independently verified literature assessment.
\item Supplied source list: Elayavalli, \emph{50 Open Problems: Ultraproduct \(\mathrm{II}_1\) Factors} (2025), Problem 48, \url{https://arxiv.org/abs/2511.20377}. The record reports access on 2026-10-16 and extraction from arXiv LaTeX; these provenance details are supplied and have not been independently verified.
\end{itemize}
\end{document}
